% ============================================================
% MATERIAL SUPLEMENTAR (versão em português, para leitura/revisão)
% Viés Geográfico e Linguístico em LLMs para Política de Poluição do Ar
% Um benchmark de 25 países, 14 modelos, 4 idiomas e 5 tarefas
% ============================================================
% Tradução fiel do supplement.tem em inglês: prosa e legendas em PT-BR;
% TODOS os números, estatísticas e valores de tabela são idênticos ao original.
% ============================================================

\documentclass[11pt,a4paper]{article}

\usepackage[utf8]{inputenc}
\usepackage[T1]{fontenc}
\usepackage[brazil]{babel}
\usepackage[margin=2.5cm]{geometry}
\usepackage{setspace}\onehalfspacing
\usepackage{amsmath,amssymb}
\usepackage{graphicx}
\usepackage{xcolor}
\usepackage{booktabs}
\usepackage{longtable}
\usepackage{array}
\usepackage{caption}\captionsetup{font=small,labelfont=bf}
\usepackage[authoryear,round]{natbib}
\bibliographystyle{apalike}
\usepackage[colorlinks=true,linkcolor=blue!50!black,citecolor=blue!50!black,urlcolor=blue!50!black]{hyperref}

% Numerar floats, equações e seções com prefixo "S".
\renewcommand{\thesection}{S\arabic{section}}
\renewcommand{\thetable}{S\arabic{table}}
\renewcommand{\thefigure}{S\arabic{figure}}
\captionsetup{labelsep=period}

\title{Material Suplementar\\[0.3em]
\large De qual regulação o modelo sabe? Confiabilidade da informação regulatória\\
recuperada por LLMs para a gestão ambiental pública em 25 países}
\author{}
\date{}

% AUTO-GERADO por code/analysis/make_numbers_tex.py — NAO EDITAR. Rodou o pipeline, este arquivo muda e o texto acompanha.
\newcommand{\naccgn}{0.572}
\newcommand{\naccgs}{0.521}
\newcommand{\nalpha}{0.527}
\newcommand{\nbayes}{-0.049}
\newcommand{\nbayeshi}{-0.006}
\newcommand{\nbayeslo}{-0.088}
\newcommand{\nbayesp}{0.987}
\newcommand{\nbgdt}{0.30}
\newcommand{\nbytesp}{0.49}
\newcommand{\nbytespartial}{+0.125}
\newcommand{\nbytespartialp}{0.56}
\newcommand{\nbytesrho}{+0.146}
\newcommand{\ncabra}{0.351}
\newcommand{\ncells}{8{,}244}
\newcommand{\ncellsen}{6{,}573}
\newcommand{\ncellsenpt}{6.573}
\newcommand{\ncellspt}{8.244}
\newcommand{\ncodecells}{6{,}434}
\newcommand{\ncodecellspt}{6.434}
\newcommand{\ncodecovtone}{99.8}
\newcommand{\ncodecovttwo}{56.3}
\newcommand{\ncodecovtthree}{71.8}
\newcommand{\ncodegap}{+14.7}
\newcommand{\ncodegapci}{[+5.6,+23.9]}
\newcommand{\ncodegapp}{0.012}
\newcommand{\ncodegn}{0.61}
\newcommand{\ncodegs}{0.46}
\newcommand{\ncodeor}{0.53}
\newcommand{\ncodeorci}{[0.48,0.58]}
\newcommand{\ndeepseek}{0.687}
\newcommand{\ndid}{+0.5}
\newcommand{\ndidmanager}{+5.3}
\newcommand{\ndidneutral}{+4.8}
\newcommand{\ndidtwo}{+0.48}
\newcommand{\ndivergent}{-5.2}
\newcommand{\ndivergentn}{298}
\newcommand{\ndupchange}{3.1}
\newcommand{\ndupmean}{0.0002}
\newcommand{\ndupresp}{21.9}
\newcommand{\ndupscore}{11.5}
\newcommand{\nevalue}{1.71}
\newcommand{\nevaluelim}{1.31}
\newcommand{\nfaithful}{-4.8}
\newcommand{\nfaithfuln}{541}
\newcommand{\nfaithfulp}{2\times10^{-8}}
\newcommand{\nfidelityp}{0.80}
\newcommand{\nfloordelta}{-0.41}
\newcommand{\nfloorequal}{-0.383}
\newcommand{\nfloorpca}{-0.435}
\newcommand{\nfloorprimary}{-0.410}
\newcommand{\nfloorrecall}{0.434}
\newcommand{\nfloorrest}{0.612}
\newcommand{\nfloorrestshort}{0.61}
\newcommand{\ngapequal}{+4.2}
\newcommand{\ngapequaltwo}{+4.22}
\newcommand{\ngapfactual}{+8.9}
\newcommand{\ngappca}{+4.7}
\newcommand{\ngappcatwo}{+4.72}
\newcommand{\ngapprimarytwo}{+5.05}
\newcommand{\nhtwocmean}{-5.0}
\newcommand{\nhtwocneg}{8}
\newcommand{\nhtwocode}{-3.8}
\newcommand{\nhtwocodeall}{-5.0}
\newcommand{\nhtwocodealln}{318}
\newcommand{\nhtwocodeallp}{0.048}
\newcommand{\nhtwocodees}{-2.1}
\newcommand{\nhtwocodeesn}{165}
\newcommand{\nhtwocoden}{333}
\newcommand{\nhtwocodep}{0.004}
\newcommand{\nhtwocodept}{-9.1}
\newcommand{\nhtwocodeptn}{132}
\newcommand{\nhtwocodeptp}{0.019}
\newcommand{\nhtwocse}{1.1}
\newcommand{\nhtwoct}{-4.6}
\newcommand{\nhtwoctp}{0.002}
\newcommand{\nhtwomixedcountry}{-5.0}
\newcommand{\nhtwomixedcountryp}{6\times10^{-7}}
\newcommand{\nhtwomixedcountryse}{1.0}
\newcommand{\nhtwomixedmodel}{-4.8}
\newcommand{\nhtwomixedmodelp}{2\times10^{-9}}
\newcommand{\nhtwomneg}{12}
\newcommand{\nhtwomwp}{0.058}
\newcommand{\nhtwoorig}{-1.1}
\newcommand{\nhtwoorign}{192}
\newcommand{\nhtwoorigp}{0.44}
\newcommand{\nhtwopanel}{-8.6}
\newcommand{\nhtwopaneln}{314}
\newcommand{\nhtwopanelp}{4\times10^{-22}}
\newcommand{\nhfive}{+12.8}
\newcommand{\nhfouranglo}{+0.024}
\newcommand{\nhfouranglop}{0.003}
\newcommand{\nhfourbeta}{+0.031}
\newcommand{\nhfourcob}{+0.030}
\newcommand{\nhfourcobp}{1\times10^{-9}}
\newcommand{\nhfourdeficit}{0.196}
\newcommand{\nhfourdep}{0.462}
\newcommand{\nhfourind}{0.658}
\newcommand{\nhfourjcob}{+0.013}
\newcommand{\nhfourjcobp}{0.057}
\newcommand{\nhfourjhdi}{+0.027}
\newcommand{\nhfourjhdip}{5\times10^{-5}}
\newcommand{\nhfourlang}{+0.004}
\newcommand{\nhfourlangp}{0.41}
\newcommand{\nhfourn}{6{,}573}
\newcommand{\nhfournotfive}{+0.026}
\newcommand{\nhfournotfivep}{3\times10^{-5}}
\newcommand{\nhfournotfour}{+0.034}
\newcommand{\nhfournotfourp}{6\times10^{-10}}
\newcommand{\nhfourp}{6\times10^{-10}}
\newcommand{\nhfourpartial}{+0.19}
\newcommand{\nhfourpartialp}{0.37}
\newcommand{\nhfourpartialthree}{+0.190}
\newcommand{\nhfourplacebo}{-0.031}
\newcommand{\nhfourse}{0.005}
\newcommand{\nhfoursite}{+0.37}
\newcommand{\nhfoursitep}{0.070}
\newcommand{\nhfoursitethree}{+0.369}
\newcommand{\nhfourtonevsfive}{+0.046}
\newcommand{\nhfourtonevsfivep}{1\times10^{-11}}
\newcommand{\nhfourwiki}{+0.05}
\newcommand{\nhfourwikip}{0.82}
\newcommand{\nhfourwikithree}{+0.048}
\newcommand{\nhthreebra}{0.258}
\newcommand{\nhthreebradelta}{-0.26}
\newcommand{\nhthreebrallama}{0.313}
\newcommand{\nhthreeequal}{-0.507}
\newcommand{\nhthreelus}{0.304}
\newcommand{\nhthreelusdelta}{-0.17}
\newcommand{\nhthreeluslama}{0.362}
\newcommand{\nhthreepca}{-0.435}
\newcommand{\nhthreeprimary}{-0.476}
\newcommand{\nicc}{0.558}
\newcommand{\niccpanel}{0.791}
\newcommand{\nincorrectn}{759}
\newcommand{\nindia}{0.651}
\newcommand{\njtgngpt}{0.38}
\newcommand{\njtgnhaiku}{0.44}
\newcommand{\njtgsgpt}{0.21}
\newcommand{\njtgshaiku}{0.37}
\newcommand{\njudgeclaude}{0.451}
\newcommand{\njudgedeepseek}{0.379}
\newcommand{\njudgegemini}{0.617}
\newcommand{\nladdergn}{0.77}
\newcommand{\nladdergs}{0.42}
\newcommand{\nladderhit}{281}
\newcommand{\nllamab}{0.438}
\newcommand{\nloocmax}{-5.5}
\newcommand{\nloocmin}{-4.3}
\newcommand{\nloogapmax}{+6.0}
\newcommand{\nloogapmin}{+4.4}
\newcommand{\nloommax}{-5.2}
\newcommand{\nloommin}{-4.5}
\newcommand{\nloorhomax}{+0.5}
\newcommand{\nloorhomaxabs}{0.48}
\newcommand{\nloorhomin}{+0.3}
\newcommand{\nmanip}{50.2}
\newcommand{\nmanipn}{2183}
\newcommand{\nmde}{0.56}
\newcommand{\nmedab}{+0.67}
\newcommand{\nmedabx}{+0.18}
\newcommand{\nmedb}{+0.27}
\newcommand{\nmedc}{+0.31}
\newcommand{\nmedhi}{+0.59}
\newcommand{\nmedlo}{-0.27}
\newcommand{\nmixedbeta}{-0.063}
\newcommand{\nmixedp}{0.086}
\newcommand{\nmixedse}{0.036}
\newcommand{\nnativa}{5.0}
\newcommand{\nnativaequal}{-4.8}
\newcommand{\nnativaes}{-4.5}
\newcommand{\nnativaesp}{2\times10^{-10}}
\newcommand{\nnativafactual}{-6.1}
\newcommand{\nnativahi}{-11.4}
\newcommand{\nnativahip}{4\times10^{-4}}
\newcommand{\nnativap}{5\times10^{-16}}
\newcommand{\nnativapt}{-4.2}
\newcommand{\nnativaptp}{8\times10^{-5}}
\newcommand{\nnativasigned}{-5.0}
\newcommand{\nnpares}{839}
\newcommand{\nnpareses}{457}
\newcommand{\nnpareshi}{71}
\newcommand{\nnparespt}{311}
\newcommand{\nnsabst}{33}
\newcommand{\nnscells}{256}
\newcommand{\nnsdk}{31}
\newcommand{\nnsfab}{168}
\newcommand{\nntasktone}{1{,}279}
\newcommand{\nntaskttwo}{1{,}335}
\newcommand{\nntasktthree}{1{,}329}
\newcommand{\nntasktfour}{1{,}310}
\newcommand{\nntasktfive}{1{,}320}
\newcommand{\norcodetone}{0.32}
\newcommand{\norcodettwo}{0.51}
\newcommand{\norcodetthree}{0.65}
\newcommand{\nort}{0.32}
\newcommand{\nortone}{0.32}
\newcommand{\nortoneci}{[0.27,0.39]}
\newcommand{\nortonegn}{50.7}
\newcommand{\nortonegs}{30.9}
\newcommand{\norttwo}{0.47}
\newcommand{\norttwoci}{[0.39,0.55]}
\newcommand{\norttwogn}{50.7}
\newcommand{\norttwogs}{37.8}
\newcommand{\nortthree}{0.69}
\newcommand{\nortthreeci}{[0.59,0.80]}
\newcommand{\nortthreegn}{59.1}
\newcommand{\nortthreegs}{52.3}
\newcommand{\nortfour}{0.33}
\newcommand{\nortfourci}{[0.26,0.40]}
\newcommand{\nortfourgn}{62.0}
\newcommand{\nortfourgs}{44.7}
\newcommand{\nortfive}{3.58}
\newcommand{\nortfiveci}{[0.84,15.17]}
\newcommand{\nortfivegn}{99.6}
\newcommand{\nortfivegs}{99.7}
\newcommand{\nortci}{[0.27,0.39]}
\newcommand{\nortexcl}{0.39}
\newcommand{\nortexclci}{[0.32,0.49]}
\newcommand{\nortgn}{50.7}
\newcommand{\nortgs}{31.0}
\newcommand{\nortladder}{0.14}
\newcommand{\nortladderci}{[0.11,0.16]}
\newcommand{\nortpre}{0.27}
\newcommand{\nortpreci}{[0.21,0.34]}
\newcommand{\npersonap}{0.29}
\newcommand{\npersonapthree}{0.287}
\newcommand{\npower}{77}
\newcommand{\npowerweak}{48}
\newcommand{\nprehfour}{+0.381}
\newcommand{\nprehfourp}{0.18}
\newcommand{\nprejoshi}{+0.131}
\newcommand{\nprejoship}{0.64}
\newcommand{\npremkp}{0.92}
\newcommand{\npremks}{3}
\newcommand{\npremkz}{+0.10}
\newcommand{\nprerho}{+0.079}
\newcommand{\nprerhop}{0.78}
\newcommand{\nregdelta}{-0.48}
\newcommand{\nrepdiff}{8.5}
\newcommand{\nrepsd}{10.1}
\newcommand{\nrestmean}{0.557}
\newcommand{\nrhat}{1.001}
\newcommand{\nrho}{0.36}
\newcommand{\nrhocobhdi}{0.65}
\newcommand{\nrhoequal}{0.35}
\newcommand{\nrhoequalthree}{+0.355}
\newcommand{\nrhomkp}{0.065}
\newcommand{\nrhop}{0.076}
\newcommand{\nrhopca}{0.38}
\newcommand{\nrhopcathree}{+0.380}
\newcommand{\nrhopre}{0.08}
\newcommand{\nrhoprimarythree}{+0.361}
\newcommand{\nrhothree}{+0.361}
\newcommand{\nsitepartialp}{0.37}
\newcommand{\nsitepthree}{0.070}
\newcommand{\nstmtp}{0.23}
\newcommand{\nstmtpartial}{+0.125}
\newcommand{\nstmtpartialp}{0.56}
\newcommand{\nstmtrho}{+0.248}
\newcommand{\nstrictgn}{0.51}
\newcommand{\nstrictgs}{0.31}
\newcommand{\ntasktone}{0.393}
\newcommand{\ntasktoneshort}{0.39}
\newcommand{\ntaskttwo}{0.473}
\newcommand{\ntasktthree}{0.548}
\newcommand{\ntasktfour}{0.520}
\newcommand{\ntasktfive}{0.769}
\newcommand{\ntaxlow}{+1.9}
\newcommand{\ntaxlowp}{0.482}
\newcommand{\ntaxlowtwo}{+1.94}
\newcommand{\ntaxmed}{+2.1}
\newcommand{\ntaxmedp}{0.329}
\newcommand{\ntaxmedtwo}{+2.14}
\newcommand{\ntaxunctad}{+5.05}
\newcommand{\ntaxvh}{+2.8}
\newcommand{\ntaxvhp}{0.206}
\newcommand{\ntaxvhtwo}{+2.78}
\newcommand{\ntiergap}{+5.0}
\newcommand{\ntiergapci}{[+1.6,+8.4]}
\newcommand{\ntiergapperm}{0.019}
\newcommand{\ntiergappermpct}{1.9}
\newcommand{\ntiergappre}{+6.1}
\newcommand{\ntiergappreci}{[+1.8,+10.1]}
\newcommand{\ntiergapprep}{0.12}
\newcommand{\ntiergapwave}{+5.9}
\newcommand{\ntiergapwaveci}{[+2.3,+9.6]}
\newcommand{\ntiergapwavep}{0.062}
\newcommand{\ntonecells}{1{,}617}
\newcommand{\ntonecellspt}{1.617}
\newcommand{\ntonecodecells}{1{,}614}
\newcommand{\ntonecodecellspt}{1.614}
\newcommand{\ntrim}{-4.5}
\newcommand{\ntrimp}{6\times10^{-20}}
\newcommand{\nverifdiff}{-9.4}
\newcommand{\nverifdiffabs}{9.4}
\newcommand{\nverifdisc}{218}
\newcommand{\nverifen}{71.9}
\newcommand{\nverifes}{-13.4}
\newcommand{\nverifhi}{-39.5}
\newcommand{\nverifnat}{62.5}
\newcommand{\nverifp}{9\times10^{-4}}
\newcommand{\nverifpt}{+3.5}
\newcommand{\nverifworse}{134}

\begin{document}
\maketitle

\noindent\emph{Reprodutibilidade.} Toda tabela orientada por dados a seguir
(Tabelas~\ref{tab:s-roster}--\ref{tab:s-groundtruth}) é regenerada a partir dos
artefatos confirmatórios brutos por
\texttt{code/analysis/make\_supplement\_tables.py}; nenhum valor é digitado à mão.
Nesta versão de leitura em português as tabelas geradas são as mesmas da versão
em inglês (com legendas em inglês), para que exista uma única fonte dos números.
A pontuação é decidida por código onde o registro tem valor (T1; a parte
resolvível de T2 e T3), pela média de um painel de três fornecedores
(Gemini~2.5~Pro, Claude Sonnet~4.6, DeepSeek-V3) no resíduo de T2 e T3 e em toda a
T4, e pelo juiz da coleta em T5 (\texttt{gpt-5-mini-2025-08-07} para os quinze
países pré-especificados, \texttt{claude-haiku-4-5-20251001} para os dez da
extensão). O corpus confirmatório analisado é composto por \ncellspt{} células
pontuadas (\ncellsenpt{} com prompt em inglês; 1.671 em idioma nativo; 951 células
receberam mais de um escore em reexecuções e foram colapsadas pela média; as 56
células egípcias de T1, sem chave no registro, foram removidas) sobre 25 países e
14 deployment-stacks. Todos os insumos, o código
de pontuação, o registro de gabarito e as respostas anonimizadas serão liberados
na publicação sob CC-BY-4.0 (dados/prompts) e MIT (código).

\tableofcontents

% ============================================================
\section{Pré-especificação, desenho e seleção de países}
% ============================================================

As hipóteses, os modelos primários, o protocolo de persona, os critérios de
exclusão, as correções de comparações múltiplas e o quadro amostral de países
foram especificados antes da coleta confirmatória de dados. O desenho
pré-especificado cobre 15 países; a amostra foi subsequentemente estendida,
\emph{post registration}, para 25 países a fim de ganhar poder para o
gradiente de desenvolvimento (H1) e para os testes de mecanismo de corpus (H4) e
multiplicar os pares de idioma intra-país (H2). Ao longo do manuscrito e deste
suplemento, todo efeito calculado na amostra de 25 países é reportado ao lado de
seu valor pré-especificado de 15 países, e nenhuma conclusão pré-especificada é
revista com base apenas na extensão (Seção~\ref{sec:s-stats}).

Os países foram selecionados por amostragem teórica triangulada em três
classificações independentes: a distinção Norte/Sul Global, tomada da classificação
de economias desenvolvidas/em desenvolvimento da UNCTAD
(operacionalizando o enquadramento da colonialidade do conhecimento), a taxonomia
de seis classes de recurso linguístico de Joshi et al.\ (2020)~\citep{joshi2020state}
e os grupos de renda do Banco Mundial. Os 15 pré-especificados (Conjunto ``P'' na
Tabela~\ref{tab:s-cov-country}) abrangem quatro regiões do mundo, cinco classes de
Joshi e quatro grupos de renda; os 10 países de extensão (Conjunto ``E'')
compreendem sete referências adicionais do Norte Global (Reino Unido, Canadá,
Austrália, Coreia do Sul, França, Itália, Portugal) escolhidas para adensar a
extremidade de alto desenvolvimento do gradiente, e três países adicionais do Sul
Global com par de idioma nativo (Colômbia, Chile, Angola). A UNCTAD declara que as
economias desenvolvidas ``compreendem, em linhas gerais, a América do Norte e a
Europa, Israel, Japão, República da Coreia, Austrália e Nova Zelândia'', sendo as
demais em desenvolvimento; os 25 países caem todos no lado que essa regra lhes
atribui, Coreia do Sul inclusive. A justificativa completa do desenho está em
\texttt{docs/supplement/s1\_design\_rationale.md}.

% ============================================================
\section{Roster de deployment-stacks e configuração de serving}
% ============================================================

\paragraph{Unidade de análise.} Avaliamos cada modelo \emph{tal como servido} por
meio de uma stack de acesso fixa --- a tupla de pesos do modelo, fornecedor,
infraestrutura de serving e camada de API/endpoint --- em vez de pesos idealizados
isoladamente, porque um gestor público interage com a stack servida e qualquer
lacuna medida é uma propriedade dessa stack. A Tabela~\ref{tab:s-roster} lista as
14 stacks com suas strings exatas de API/tag e os locais de execução.

% AUTO-GENERATED by code/analysis/make_supplement_tables.py — do not edit by hand
\begin{longtable}{p{2.6cm}llp{4.4cm}lrr}
\caption{Deployment-stack roster. Each model is evaluated \emph{as served} through a fixed access stack; the API/tag string and execution venue are part of what each label denotes. Parameter counts are vendor-reported (dense or total for MoE; closed-frontier counts undisclosed). $N$ and mean are realized English-prompt confirmatory coverage.\label{tab:s-roster}}\\
\toprule
Model & Vendor & Tier & API / tag string & Venue & $N$ & Mean \\
\midrule \endfirsthead
\toprule Model & Vendor & Tier & API / tag string & Venue & $N$ & Mean \\ \midrule \endhead
DeepSeek-V3 & DeepSeek & A & \texttt{deepseek-chat} & DeepSeek API & 496 & 0.690 \\
Command R+ & Cohere & A & \texttt{command-r-plus-08-2024} & Cohere (trial) & 496 & 0.530 \\
Llama 3.3 70B & Meta & A & \texttt{llama-3.3-70b-versatile} & Groq (free) & 391 & 0.514 \\
Llama 4 Scout & Meta & A & \texttt{meta-llama/\allowbreak llama-4-scout-17b-16e-instruct} & Groq (free) & 496 & 0.500 \\
GPT-OSS 120B & OpenAI & A & \texttt{openai/\allowbreak gpt-oss-120b} & Groq (free) & 257 & 0.479 \\
Qwen3 32B & Alibaba & B & \texttt{qwen/\allowbreak qwen3-32b} & Groq (free) & 494 & 0.560 \\
Qwen3 14B & Alibaba & B & \texttt{qwen3:14b} & Ollama (local) & 496 & 0.485 \\
Phi-4 14B & Microsoft & B & \texttt{phi4} & Ollama (local) & 496 & 0.475 \\
Llama 3.1 8B & Meta & C & \texttt{llama-3.1-8b-instant} & Groq (free) & 496 & 0.437 \\
Cabra-Mistral 7B & botbot-ai & C & \texttt{cabra-mistral-7b} & Ollama (local) & 492 & 0.351 \\
Gemini 2.5 Flash & Google & D & \texttt{gemini-2.5-flash} & Google AI (free) & 486 & 0.636 \\
GPT-5-mini & OpenAI & D & \texttt{gpt-5-mini} & OpenAI API & 495 & 0.628 \\
Claude Haiku 4.5 & Anthropic & D & \texttt{claude-haiku-4-5} & Anthropic API & 486 & 0.589 \\
GPT-5 & OpenAI & E & \texttt{gpt-5} & OpenAI API & 496 & 0.676 \\
\bottomrule
\end{longtable}


\paragraph{Determinismo e tratamento de seed.} Todas as stacks foram consultadas a
temperatura~0,3 com a configuração de máximo determinismo que cada provedor
documenta, e não com uma superfície idêntica de parâmetros, porque as superfícies
diferem. O suporte a seed é heterogêneo: stacks servidas localmente (Ollama)
honram uma seed fixa; a OpenAI expõe uma seed advisory; a Messages API da
Anthropic não expõe parâmetro de seed; e, para os demais endpoints hospedados, a
seed é advisory ou não é exposta. Para cada chamada registramos a seed solicitada
e um campo \texttt{seed\_status} (\texttt{sent}, \texttt{logged-only} ou
\texttt{not-supported}) nos metadados de execução liberados, de modo que o
determinismo de fato alcançado seja auditável em vez de presumido. Duas replicações
estocásticas por célula (prompt, modelo, persona) estimam a variabilidade
intra-célula.

\paragraph{Janela de aquisição.} O comportamento do modelo servido pode derivar à
medida que fornecedores atualizam a infraestrutura, de modo que as acurácias
reportadas são estimativas point-in-time relativas à janela de coleta. A tag
fixada, a string de versão e o timestamp de aquisição são registrados para cada
chamada. A coleta confirmatória ocorreu de 2026-04-26 a 2026-06-12 (UTC),
abrangendo o piloto, a coleta pré-especificada de 15 países e a extensão
post-registration para 25 países. A Tabela~\ref{tab:s-acquisition} dá o timestamp
da primeira e da última chamada sem erro por provedor, junto com a documentação
oficial de API consultada; os timestamps completos por chamada são liberados com o
conjunto de dados.

% AUTO-GENERATED by code/analysis/make_supplement_tables.py
\begin{longtable}{lp{2.0cm}p{2.0cm}p{5.3cm}}
\caption{Per-provider acquisition window and official API documentation. Reported accuracies are point-in-time estimates relative to this window; first/last are the earliest and latest non-error call timestamps (UTC) recorded in the run logs. The window spans the pilot, the pre-registered 15-country collection, and the post-registration extension.\label{tab:s-acquisition}}\\
\toprule
Provider (venue) & First (UTC) & Last (UTC) & API documentation \\
\midrule \endfirsthead
\toprule Provider (venue) & First (UTC) & Last (UTC) & API documentation \\ \midrule \endhead
Anthropic API & 2026-04-27 & 2026-06-12 & \url{https://docs.anthropic.com/en/api/messages} \\
Cohere (trial) & 2026-04-26 & 2026-06-11 & \url{https://docs.cohere.com/reference/chat} \\
DeepSeek API & 2026-04-26 & 2026-06-11 & \url{https://api-docs.deepseek.com} \\
Google AI (free) & 2026-04-27 & 2026-06-11 & \url{https://ai.google.dev/api} \\
Groq (free) & 2026-04-26 & 2026-06-11 & \url{https://console.groq.com/docs/api-reference} \\
Ollama (local) & 2026-04-27 & 2026-06-12 & \url{https://github.com/ollama/ollama/blob/main/docs/api.md} \\
OpenAI API & 2026-04-27 & 2026-06-12 & \url{https://platform.openai.com/docs/api-reference/chat} \\
\midrule
\textbf{All providers} & 2026-04-26 & 2026-06-12 & --- \\
\bottomrule
\end{longtable}


% ============================================================
\section{Cobertura}
% ============================================================

A cobertura por célula não é uniforme ao longo do roster: endpoints gratuitos e com
limite de taxa retornaram menos respostas admissíveis do que as APIs medidas, e um
candidato (Gemini~2.5~Flash em lotes nativos iniciais) e um sexto juiz
(Command~R+) atingiram cotas dos provedores. Reportamos a cobertura realizada em
vez de descartar silenciosamente células sub-cobertas. A
Tabela~\ref{tab:s-cov-task} dá a matriz modelo~$\times$~tarefa para prompts em
inglês (os totais por coluna igualam as contagens por tarefa do texto principal);
a Tabela~\ref{tab:s-cov-country} dá as contagens por país em inglês, nativo e total
com camada e conjunto de pré-especificação; a Tabela~\ref{tab:s-native} lista os pares
de idioma nativo usados em H2. A matriz de cobertura completa
modelo~$\times$~país~$\times$~tarefa, com o motivo de cota ou exclusão de cada
déficit, está incluída no release.

% AUTO-GENERATED by code/analysis/make_supplement_tables.py
\begin{longtable}{p{2.8cm}rrrrrr}
\caption{Coverage matrix: realized English-prompt responses per model $\times$ task (T1 standard, T2 local datum, T3 health synthesis, T4 instruments, T5 recommendation). No cell is silently dropped; shortfalls are quota-driven (Methods).\label{tab:s-cov-task}}\\
\toprule
Model & T1 & T2 & T3 & T4 & T5 & Total \\
\midrule \endfirsthead
\toprule Model & T1 & T2 & T3 & T4 & T5 & Total \\ \midrule \endhead
GPT-5 & 96 & 100 & 100 & 100 & 100 & 496 \\
DeepSeek-V3 & 96 & 100 & 100 & 100 & 100 & 496 \\
Command R+ & 96 & 100 & 100 & 100 & 100 & 496 \\
Phi-4 14B & 96 & 100 & 100 & 100 & 100 & 496 \\
Qwen3 14B & 96 & 100 & 100 & 100 & 100 & 496 \\
Llama 4 Scout & 96 & 100 & 100 & 100 & 100 & 496 \\
Llama 3.1 8B & 96 & 100 & 100 & 100 & 100 & 496 \\
GPT-5-mini & 96 & 100 & 100 & 99 & 100 & 495 \\
Qwen3 32B & 96 & 100 & 100 & 99 & 99 & 494 \\
Cabra-Mistral 7B & 95 & 97 & 100 & 100 & 100 & 492 \\
Gemini 2.5 Flash & 94 & 98 & 98 & 98 & 98 & 486 \\
Claude Haiku 4.5 & 94 & 98 & 98 & 98 & 98 & 486 \\
Llama 3.3 70B & 78 & 81 & 80 & 76 & 76 & 391 \\
GPT-OSS 120B & 54 & 61 & 53 & 40 & 49 & 257 \\
\midrule
\textbf{Total} & 1279 & 1335 & 1329 & 1310 & 1320 & 6573 \\
\bottomrule
\end{longtable}

% AUTO-GENERATED by code/analysis/make_supplement_tables.py
\begin{longtable}{llcrrr}
\caption{Coverage by country: tier (GN/GS), sample set (P = pre-specified 15; E = extension), and realized responses (English, native-language, total).\label{tab:s-cov-country}}\\
\toprule
Country & Tier & Set & English & Native & Total \\
\midrule \endfirsthead
\toprule Country & Tier & Set & English & Native & Total \\ \midrule \endhead
Colombia & GS & E & 260 & 251 & 511 \\
Chile & GS & E & 260 & 241 & 501 \\
Portugal & GN & E & 261 & 240 & 501 \\
Angola & GS & E & 260 & 240 & 500 \\
Argentina & GS & P & 280 & 139 & 419 \\
India & GS & P & 279 & 138 & 417 \\
Mexico & GS & P & 269 & 140 & 409 \\
Brazil & GS & P & 260 & 141 & 401 \\
Peru & GS & P & 256 & 141 & 397 \\
Germany & GN & P & 279 & 0 & 279 \\
Indonesia & GS & P & 279 & 0 & 279 \\
Japan & GN & P & 279 & 0 & 279 \\
Kenya & GS & P & 279 & 0 & 279 \\
Bangladesh & GS & P & 278 & 0 & 278 \\
United Kingdom & GN & E & 265 & 0 & 265 \\
Canada & GN & E & 260 & 0 & 260 \\
Australia & GN & E & 260 & 0 & 260 \\
France & GN & E & 260 & 0 & 260 \\
Italy & GN & E & 260 & 0 & 260 \\
South Korea & GN & E & 258 & 0 & 258 \\
Nigeria & GS & P & 257 & 0 & 257 \\
Philippines & GS & P & 256 & 0 & 256 \\
United States & GN & P & 252 & 0 & 252 \\
South Africa & GS & P & 242 & 0 & 242 \\
Egypt & GS & P & 224 & 0 & 224 \\
\midrule
\textbf{Total (25)} & & & 6573 & 1671 & 8244 \\
\bottomrule
\end{longtable}

% AUTO-GENERATED by code/analysis/make_supplement_tables.py
\begin{longtable}{llr}
\caption{Native-language confirmatory responses by language and country (H2 within-country English/native contrast).\label{tab:s-native}}\\
\toprule Native language (Joshi) & Country & $N$ (native) \\ \midrule \endfirsthead
\toprule Native language (Joshi) & Country & $N$ (native) \\ \midrule \endhead
Portuguese (4) & Portugal & 240 \\
 & Angola & 240 \\
 & Brazil & 141 \\
\addlinespace
Spanish (5) & Colombia & 251 \\
 & Chile & 241 \\
 & Peru & 141 \\
 & Mexico & 140 \\
 & Argentina & 139 \\
\addlinespace
Hindi (4) & India & 138 \\
\addlinespace
\bottomrule
\end{longtable}


% ============================================================
\section{Covariáveis dos países}
% ============================================================

A Tabela~\ref{tab:s-covariates} reporta as covariáveis de desenvolvimento e de
representação no corpus usadas na análise de mecanismo (H4). As medidas de
corpus do idioma (classe Joshi; tamanhos mC4/OSCAR reportados no texto
principal) indexam quanto texto existe \emph{no} idioma de um país e são o canal
candidato para a penalidade de idioma nativo (H2). As medidas de
cobertura de país (comprimento em bytes do artigo na Wikipédia em inglês,
sitelinks do Wikidata, statements do Wikidata) indexam quão amplamente o corpus
enciclopédico é \emph{sobre} um país, independentemente do idioma, e são o canal
candidato para a lacuna geográfica (H1/H4). Como a lacuna geográfica é medida
\emph{em inglês}, uma lacuna Norte/Sul residual não pode ser um efeito de corpus do
idioma, razão pela qual H4 isola a cobertura de país para H1 e o tamanho do corpus
do idioma para H2. Todas as entidades do Wikidata são identificadas por um QID
resolvível e foram recuperadas dentro da janela de coleta.

% AUTO-GENERATED by code/analysis/make_supplement_tables.py
\begin{longtable}{llrrrrr}
\caption{Country covariates. HDI: UNDP HDR 2023--24 (2022 data). Joshi: linguistic-resource class of the dominant official language~\citep{joshi2020state}. Corpus measures from Wikimedia/Wikidata (entity QID resolvable): English-article byte length, Wikidata sitelinks (language editions), and Wikidata statements.\label{tab:s-covariates}}\\
\toprule
Country & Tier & HDI & Joshi & EN-wiki bytes & Sitelinks & WD stmts \\
\midrule \endfirsthead
\toprule Country & Tier & HDI & Joshi & EN-wiki bytes & Sitelinks & WD stmts \\ \midrule \endhead
Germany & GN & 0.950 & 5 & 207,810 & 403 & 1402 \\
Australia & GN & 0.946 & 5 & 233,095 & 395 & 1674 \\
United Kingdom & GN & 0.940 & 5 & 368,888 & 391 & 1085 \\
Canada & GN & 0.935 & 5 & 281,016 & 393 & 1346 \\
South Korea & GN & 0.929 & 4 & 275,291 & 350 & 746 \\
United States & GN & 0.927 & 5 & 383,914 & 424 & 1710 \\
Japan & GN & 0.920 & 5 & 205,844 & 414 & 995 \\
France & GN & 0.910 & 5 & 279,498 & 413 & 1196 \\
Italy & GN & 0.906 & 4 & 305,234 & 405 & 1039 \\
Portugal & GN & 0.874 & 4 & 324,038 & 369 & 751 \\
Chile & GS & 0.860 & 5 & 232,108 & 379 & 689 \\
Argentina & GS & 0.849 & 5 & 267,662 & 369 & 881 \\
Mexico & GS & 0.781 & 5 & 263,757 & 398 & 1286 \\
Peru & GS & 0.762 & 5 & 239,544 & 332 & 867 \\
Brazil & GS & 0.760 & 4 & 301,844 & 384 & 1151 \\
Colombia & GS & 0.758 & 5 & 303,401 & 363 & 774 \\
Egypt & GS & 0.728 & 5 & 284,547 & 387 & 757 \\
South Africa & GS & 0.717 & 5 & 270,861 & 365 & 750 \\
Indonesia & GS & 0.713 & 3 & 226,349 & 362 & 1005 \\
Philippines & GS & 0.710 & 5 & 493,981 & 331 & 963 \\
Bangladesh & GS & 0.670 & 3 & 275,786 & 343 & 814 \\
India & GS & 0.644 & 5 & 310,681 & 399 & 1679 \\
Kenya & GS & 0.601 & 5 & 216,189 & 334 & 699 \\
Angola & GS & 0.591 & 4 & 180,098 & 322 & 676 \\
Nigeria & GS & 0.548 & 5 & 266,041 & 349 & 1194 \\
\bottomrule
\end{longtable}


% ============================================================
\section{Registro de gabarito (T1)}
% ============================================================

A validade do gradiente geográfico (H1) requer que o gabarito seja igualmente
disponível e igualmente verificável entre os países. Para T1 (o padrão anual
vinculante de PM\textsubscript{2.5}) ancoramos cada país a uma fonte oficial
primária por auditoria documental: o texto-fonte foi recuperado, o valor foi lido
verbatim, e a entrada registra o valor, o status, a fonte-de-registro, uma URL
resolvível e o método de recuperação (Tabela~\ref{tab:s-groundtruth}). Três países
\emph{não têm} padrão nacional de PM\textsubscript{2.5}, confirmado pelo texto
oficial (a NESREA da Nigéria regula apenas PM\textsubscript{10}; a Argentina não
tem valor federal vinculante; a Angola não tem legislação nacional de qualidade do
ar); nesses casos, um modelo fiel deve recusar-se a indicar um limiar inexistente
em vez de fabricá-lo. Essa proveniência é mais reprodutível do que a asserção de
especialistas: qualquer leitor pode reverificar cada valor a partir da fonte
oficial citada. Nenhuma referência no registro é não-verificada, e nenhum valor foi
transposto de fontes secundárias sem confirmação contra o texto primário.

% AUTO-GENERATED by code/analysis/make_supplement_tables.py
\footnotesize
\begin{longtable}{p{1.9cm}p{4.3cm}p{1.9cm}p{4.6cm}}
\caption{Ground-truth registry for T1 (binding annual PM\textsubscript{2.5} standard), all 25 countries, anchored to official primary sources by documentary audit. All entries are of source type \emph{official primary} unless noted. Status \emph{verified} = value read from the official text; \emph{no standard} = official confirmation that no national standard exists; \emph{partial} = the value could not be confirmed in an accessible official source (Bangladesh: a 2022 revision is known; Egypt: no value extracted, excluded from T1). Full URLs and verbatim excerpts are in the released registry JSONL.\label{tab:s-groundtruth}}\\
\toprule
Country & Value & Status & Official source \\
\midrule \endfirsthead
\toprule Country & Value & Status & Official source \\ \midrule \endhead
Angola & NO national air-quality legislation / standard & no standard & No national ambient air quality regulation (World Bank/WHO note) \\
Argentina & NO single binding federal PM\textsubscript{2.5} standard (Ley 20.284/1973 framework; PM\textsubscript{2.5} limits set subnationally) & no standard & Ley 20.284/1973 \\
Australia & 8 $\mu$g/m\textsuperscript{3} (annual; 7 target 2025) & verified & National Environment Protection (Ambient Air Quality) Measure \\
Bangladesh & 35 $\mu$g/m\textsuperscript{3} (annual; 24-h 65) per Air Pollution (Control) Rules 2022, S.R.O. 255-Law/2022; superseded 15 (S.R.O. 220-Law/2005) & verified & Air Pollution (Control) Rules 2022 (S.R.O. 255-Law/2022), Schedule of ambient standards; confirmed in DoE, Ambient Air Quality in Bangladesh (2024), Table 3 \\
Brazil & 17 $\mu$g/m\textsuperscript{3} (annual, PI-2) & verified & CONAMA 506/2024 Annex I, DOU 09/07/2024 Ed.130 p.133 \\
Canada & 8.8 $\mu$g/m\textsuperscript{3} (annual, CAAQS 2020) & verified & Canadian Ambient Air Quality Standards (CCME) \\
Chile & 20 $\mu$g/m\textsuperscript{3} (annual) & verified & D.S. 12/2011 MMA Chile \\
Colombia & 25 $\mu$g/m\textsuperscript{3} (annual, in force since 2018; 15 from 2030; 24h=37) & verified & Resolucion 2254/2017 MinAmbiente Tabla 1 \& 2 (verbatim official scan) \\
Egypt & PM\textsubscript{2.5} limit set under Law 4/1994 (amended 9/2009, 105/2015) Executive Regulations (official text in Arabic; value not extracted from an accessible official source) & partial & Environmental Law 4/1994 (EEAA) \\
France & 25 $\mu$g/m\textsuperscript{3} (annual, in force; 10 from 2030) & verified & EU Ambient Air Quality Directive 2008/50/EC; (EU) 2024/2881 \\
Germany & 25 $\mu$g/m\textsuperscript{3} (annual, in force 2026; 10 from 2030) & verified & EU Ambient Air Quality Directive 2008/50/EC; revised (EU) 2024/2881 \\
India & 40 $\mu$g/m\textsuperscript{3} (annual) & verified & CPCB National Ambient Air Quality Standards (2009) \\
Indonesia & 15 $\mu$g/m\textsuperscript{3} (annual) & verified & PP No. 22/2021 Lampiran VII (Indonesia) \\
Italy & 25 $\mu$g/m\textsuperscript{3} (annual, in force; 10 from 2030) & verified & EU Ambient Air Quality Directive 2008/50/EC; (EU) 2024/2881 \\
Japan & 15 $\mu$g/m\textsuperscript{3} (annual) & verified & Ministry of the Environment Japan, Environmental Quality Standards (2009) \\
Kenya & 35 $\mu$g/m\textsuperscript{3} (annual; 24h 75) & verified & First Schedule, EMCA (Air Quality) Regulations 2014 (Kenya Gazette Vol. CXVI No.46) \\
Mexico & 10 $\mu$g/m\textsuperscript{3} (annual) & verified & NOM-025-SSA1-2021, Diario Oficial de la Federacion 27/10/2021 \\
Nigeria & NO national PM\textsubscript{2.5} standard (PM10 annual=60 $\mu$g/m\textsuperscript{3}, 24h=150) & no standard & NESREA National Environmental (Air Quality Control) Regulations, Schedule XIII (Ambient Air Quality Standards) \\
Peru & 25 $\mu$g/m\textsuperscript{3} (annual) & verified & D.S. 003-2017-MINAM (Peru) \\
Philippines & 25 $\mu$g/m\textsuperscript{3} (annual, 1-year guideline) & verified & DENR DAO (RA 8749 Clean Air Act); 25 ug/Nm3 annual / 50 ug/Nm3 24h \\
Portugal & 25 $\mu$g/m\textsuperscript{3} (annual, in force; 10 from 2030) & verified & EU Ambient Air Quality Directive 2008/50/EC; (EU) 2024/2881 \\
South Africa & 20 $\mu$g/m\textsuperscript{3} (annual) & verified & SA NAAQS under NEMAQA (Act 39/2004), 2012 \\
South Korea & 15 $\mu$g/m\textsuperscript{3} (annual) & verified & Ministry of Environment, Korea \\
United Kingdom & 20 $\mu$g/m\textsuperscript{3} (annual, limit value in force since 1 Jan 2020; England target 10 by 2040) & verified & Air Quality Standards Regulations 2010, Sch. 2 (Stage 2 limit value); DEFRA National air quality objectives table, 2023-04-03 \\
United States & 9.0 $\mu$g/m\textsuperscript{3} (annual, primary NAAQS, 2024) & verified & EPA Final Reconsideration of PM NAAQS, Federal Register 2024-02637 \\
\bottomrule
\end{longtable}
\normalsize


% ============================================================
\section{Resultado composto e rubrica de pontuação}
% ============================================================

Cada resposta foi pontuada pelo juiz primário contra a entrada do registro de
gabarito (T1, T2, T4) ou a rubrica validada (T3, T5), produzindo cinco
subcomponentes, cada um normalizado a $[0,1]$:
\begin{enumerate}\setlength{\itemsep}{1pt}
  \item Acurácia factual contra o valor do registro (binária para T1;
        crédito parcial para T2--T4).
  \item Completude contextual pontuada contra uma rubrica pré-validada
        (T3, T4).
  \item Qualidade de citação, verificada por checagem de DOI/fonte oficial.
  \item Calibração, o alinhamento entre a confiança declarada e a correção.
  \item Ausência de alucinação, codificada de forma binária ($0$ alucinado,
        $1$ fiel).
\end{enumerate}
O composto primário usa pesos iguais. Duas ponderações de sensibilidade
pré-especificadas são reportadas --- um conjunto derivado por PCA estimado no
piloto estendido do Brasil, e um conjunto especificado pelos autores (factual~0,30,
contextual~0,25, citação~0,15, alucinação~0,15, aplicada~0,15) --- e uma conclusão
é tratada como robusta apenas se valer sob as três. O composto é calculado por
célula (país, modelo, persona) para suportar o teste de
diferença-em-diferenças de H6. Cada registro pontuado guarda também a justificativa
textual do juiz; por exemplo, o juiz \texttt{gpt-5-mini} pontuou uma resposta do
Llama~4~Scout sobre o padrão T1 do Brasil em $0,15$ com a justificativa de que ela
``gives wrong regulation (491/2018 vs correct 506/2024) and wrong current
value\ldots\ included a specific but incorrect citation and fabricated timing'',
ilustrando como os subcomponentes de acurácia factual e de alucinação discriminam
falha substantiva.

% ============================================================
\section{Confiabilidade do painel de juízes}\label{sec:s-reliability}
% ============================================================

A confiabilidade é computada sobre cada item que o painel pontuou --- os mesmos
itens de que os efeitos relatados são computados ---, e não sobre um subconjunto
estratificado de calibração. Todas as 3.190 respostas que exigiram juiz foram
pontuadas independentemente pelos três membros do painel, escolhidos por
diversidade de fornecedor para que a correlação intrafornecedor não infle a
concordância aparente. O painel substitui o desenho de juiz único do protocolo
original, no qual o juiz da coleta (\texttt{gpt-5-mini} para os quinze países
pré-especificados, \texttt{claude-haiku-4-5} para os dez da extensão) estava ele
próprio entre os modelos avaliados e, como a extensão contém sete dos dez países
do Norte Global, estava confundido com o tier.

\begin{table}[h]\centering
\caption{Confiabilidade entre juízes sobre os 3.190 itens pontuados pelo painel. A
quantidade operativa é a confiabilidade da \emph{média} do painel, ICC$(2,3)$;
pela relação de Spearman--Brown ela é alta mesmo quando juízes isolados concordam
apenas moderadamente. Fonte: \texttt{panel\_reliability\_full.json}.}
\label{tab:s-confiabilidade}
\small
\begin{tabular}{lr}
\toprule
Quantidade & Valor \\
\midrule
Juízes ($k$) & 3 (Google, Anthropic, DeepSeek) \\
Itens pontuados por todos os juízes & 3.190 \\
$\alpha$ de Krippendorff (intervalar) & 0.527 \\
ICC$(2,1)$ de juiz único & 0.558 \\
ICC$(2,3)$ da média do painel & 0.791 \\
\midrule
Pearson par a par: Claude Sonnet $\,\vert\,$ DeepSeek-V3 & 0.801 \\
\quad Gemini 2.5 Pro $\,\vert\,$ Claude Sonnet & 0.663 \\
\quad Gemini 2.5 Pro $\,\vert\,$ DeepSeek-V3 & 0.654 \\
\midrule
Composto médio por juiz: Gemini 2.5 Pro & 0.617 \\
\quad Claude Sonnet 4.6 & 0.451 \\
\quad DeepSeek-V3 & 0.379 \\
\bottomrule
\end{tabular}
\end{table}

A concordância par a par é mais forte entre Claude e DeepSeek e menor para o
Gemini, sistematicamente mais brando. Como todo efeito relatado é um contraste
entre grupos (lacuna de camada, penalidade de idioma, piso de tarefa), o
deslocamento de brandura de um juiz se cancela e não ameaça esses contrastes ---
que é precisamente a razão de a análise usar a \emph{média} do painel em vez de
qualquer juiz isolado, cuja confiabilidade sozinha seria ICC$(2,1)=0.56$.

\section{Métodos estatísticos e saídas completas dos testes confirmatórios}\label{sec:s-stats}
% ============================================================

Os três testes confirmatórios primários formam uma única família (F1) corrigida em
conjunto pelo procedimento step-down de Bonferroni--Holm ($\alpha=0,05$). Os testes
confirmatórios secundários são reportados sem ajuste dentro da família; os testes
declarados exploratórios usam FDR de Benjamini--Hochberg a $q=0,10$. As
correlações de postos usam Spearman $\rho$; a robustez de tendência monotônica usa
o teste de Mann--Kendall; o efeito de persona usa um teste de permutação em nível
de país (5{,}000 permutações). Todas as saídas a seguir são regeneradas por
\texttt{code/analysis/formal\_tests.py}, \texttt{robust\_tests.py} e
\texttt{h4\_corpus\_mechanism.py}.

\begin{table}[h]\centering
\caption{Família primária (Bonferroni--Holm): valores pré-especificados de 15 países
ao lado da extensão post-registration de 25 países. H1-por-IDH é o único teste
primário que rejeita sua hipótese nula a $n=25$; permanece abaixo do limiar de
tamanho de efeito pré-especificado $\rho\geq 0,55$ em ambas as amostras.}
\label{tab:s-primary}
\small
\begin{tabular}{lll}
\toprule
Teste & 15 países (pré-especificado) & 25 países (extensão) \\
\midrule
H1 Spearman $\rho$(acc, IDH)        & $\nprerho$ ($p=\nprerhop$), ns            & $\nrhothree$ ($p=\nrhop$), ns \\
H1 Mann--Kendall (por IDH)          & $S=\npremks$, $Z=\npremkz$, $p=\npremkp$, ns     & $p=\nrhomkp$, ns \\
H1 Spearman $\rho$(acc, Joshi)      & $\nprejoshi$ ($p=\nprejoship$), ns                      & $-0.061$, ns (desenvolvimento, não linguístico) \\
H1 vs limiar $\rho\geq 0.55$        & não atingido                       & não atingido \\
H4 $\rho$ parcial (acc, cobertura$\,\vert\,$IDH) & $\nprehfour$ ($p=\nprehfourp$), ns & $\nhfourpartialthree$ ($p=\nsitepartialp$), ns \\
H6 DiD persona (SG$-$NG)            & não recomputada em $n=15$          & $\ndidtwo$~pp, perm $p=\npersonapthree$, ns \\
Lacuna de camada (NG$-$SG), composto & $\ntiergappre$~pp, $p=\ntiergapprep$ (3 NG)        & $\ntiergap$~pp, $p=\ntiergapperm$ \\
\bottomrule
\end{tabular}
\end{table}

\begin{table}[h]\centering
\caption{Mecanismo de corpus H4, multi-medida (25 países). A amplitude de cobertura
de país (sitelinks do Wikidata) prediz a acurácia no nível de ordem zero e
sobrevive à correção de multiplicidade entre os três proxies de cobertura, mas
atenua-se até a não-significância após ajuste para desenvolvimento (IDH);
reportamos, portanto, como sugestivo/exploratório, não como mecanismo estabelecido.
O proxy pré-especificado de edição de idioma é nulo.}
\label{tab:s-h4}
\small
\begin{tabular}{llrr}
\toprule
Medida & Família & $\rho$(acc, log medida) & parcial $\,\vert\,$ IDH \\
\midrule
Bytes Wikipédia-inglês       & país     & $\nbytesrho$ ($p=\nbytesp$) & $\nbytespartial$ ($p=\nbytespartialp$) \\
Sitelinks Wikidata  & país     & $\nhfoursitethree$ ($p=\nsitepthree$) & $\nhfourpartialthree$ ($p=\nsitepartialp$) \\
Statements Wikidata          & país     & $\nstmtrho$ ($p=\nstmtp$) & $\nstmtpartial$ ($p=\nstmtpartialp$) \\
Contagem de edições Wikipédia & idioma  & $\nhfourwikithree$, ns        & ns \\
\bottomrule
\end{tabular}
\end{table}

\begin{table}[h]\centering
\caption{Achados secundários robustos (25 países), com verificações de
leave-one-country-out (LOCO) e de restrição de amplitude.}
\label{tab:s-robust}
\footnotesize
\begin{tabular}{>{\raggedright\arraybackslash}p{0.30\linewidth}>{\raggedright\arraybackslash}p{0.64\linewidth}}
\toprule
Achado & Estatística \\
\midrule
Lacuna de camada (NG$-$SG)         & $\ntiergap$~pp, IC 95\% $\ntiergapci$; permutação $p=\ntiergapperm$ \\
Lacuna de camada, faixa LOCO       & $[\nloogapmin,\nloogapmax]$~pp, estável \\
Lacuna de camada, só onda de extensão & $\ntiergapwave$~pp $\ntiergapwaveci$ (7 NG vs 3 SG), perm $p=\ntiergapwavep$ \\
Lacuna de camada, só desfechos de código (T1; T2, T3 resolvidas) & NG \ncodegn{} vs SG \ncodegs{}, $\ncodegap$~pp $\ncodegapci$, perm $p=\ncodegapp$; RC agregada \ncodeor{} $\ncodeorci$ (T1 \norcodetone, T2 \norcodettwo, T3 \norcodetthree) \\
H1 \emph{leave-one-country-out} & $\rho\in[\nloorhomin,\nloorhomax]$: nunca atinge o limiar pré-especificado de 0.55 \\
Piso de tarefa (T1{+}T2 vs T3--T5) & $\nfloorrecall$ vs $\nfloorrest$, $\delta$ de Cliff $=\nfloordelta$ \\
H3 Cabra-Mistral vs restante       & $\ncabra$ vs $0.557$, $\delta$ de Cliff $=\nregdelta$ \\
H3 pré-especificada (Cabra vs Llama 3.1 8B, prompts em português) & Brasil $\nhthreebra$ vs $\nhthreebrallama$, $\delta=\nhthreebradelta$ ($n=20$ cada); BRA+PRT+AGO $\nhthreelus$ vs $\nhthreeluslama$, $\delta=\nhthreelusdelta$ ($n=60$ cada) \\
H2 nativo vs inglês (todos)        & $\Delta=\nnativasigned$~pp, Wilcoxon $p=\nnativap$ ($n=\nnpares$); modelo misto com intercepto aleatório por país $\nhtwomixedcountry$~pp ($p=\nhtwomixedcountryp$), por modelo $\nhtwomixedmodel$~pp ($p=\nhtwomixedmodelp$); país como unidade \nhtwocneg/9 negativos, $t=\nhtwoct$, $p=\nhtwoctp$; modelo como unidade \nhtwomneg/14, Wilcoxon $p=\nhtwomwp$ \\
\quad Espanhol / Português / Hindi & $\nnativaes$ / $\nnativapt$ / $\nnativahi$~pp, todos significativos \\
\quad por instrumento de pontuação & código $\nhtwocode$ ($n=\nhtwocoden$), painel $\nhtwopanel$ ($n=\nhtwopaneln$), juiz da coleta $\nhtwoorig$~pp ns ($n=\nhtwoorign$, T5) \\
\quad só veredito de código, por idioma & português $\nhtwocodept$ ($n=\nhtwocodeptn$, $p=\nhtwocodeptp$), espanhol $\nhtwocodees$ ($n=\nhtwocodeesn$, ns), todos $\nhtwocodeall$~pp ($n=\nhtwocodealln$, $p=\nhtwocodeallp$) \\
H5 aberto vs fechado-acessível     & vantagem fechada de $\nhfive$~pp (descritivo) \\
\bottomrule
\end{tabular}
\end{table}

% ============================================================
\section{Análises adicionais de robustez e corroboração}

\emph{Sensibilidade da ponderação do composto.} Recomputamos cada efeito
principal sob a ponderação dos autores, sob pesos iguais e sob ponderação derivada
por ACP. Os dois efeitos sustentados se mantêm nas três, e o gradiente fica
abaixo do 0.55 pré-especificado nas três. Sob a ponderação primária a lacuna de
camada é de $\ngapprimarytwo$~pp e o gradiente $\rho=\nrhoprimarythree$; sob pesos iguais, $\ngapequaltwo$~pp
e $\rho=\nrhoequalthree$; sob ACP, $\ngappcatwo$~pp e $\rho=\nrhopcathree$. Um modelo misto gaussiano
com interceptos aleatórios cruzados dá $\beta=\nmixedbeta$ ($p=\nmixedp$) para o Sul
Global; a reestimação bayesiana dá média posterior $\nbayes$, HDI de 94\%
$[\nbayeslo,\nbayeshi]$, $P(\beta<0)=\nbayesp$; o E-value é \nevalue{} no ponto e \nevaluelim{} no
limite do intervalo mais próximo do nulo.

\emph{Poder sobre o nulo do gradiente.} Simulando amostras de 25 países em
correlações conhecidas, o desenho tem 77\% de poder para detectar $\rho=0.55$ e
efeito mínimo detectável de $\rho=0.56$ a 80\% de poder. O nulo exclui, portanto,
um gradiente da magnitude pré-especificada, mas não um gradiente fraco (poder de
48\% em $\rho=0.40$).

\emph{Dependência da taxonomia de camada.} A lacuna é significativa sob a
partição desenvolvidas/em desenvolvimento da UNCTAD (permutação $p=\ntiergapperm$) e não
sob nenhuma partição dos mesmos 25 países traçada apenas por IDH (mediana
$\ntaxmedtwo$~pp, $p=\ntaxmedp$; IDH $<0.800$ $\ntaxvhtwo$~pp, $p=\ntaxvhp$; IDH $<0.700$
$\ntaxlowtwo$~pp, $p=\ntaxlowp$). Relatamos isso
em vez de descansar sobre a partição favorável.

\emph{Duplicatas, em dois níveis.} No nível da resposta, 21,9\% das chaves
(modelo, prompt, réplica) têm mais de uma resposta armazenada; a que entra na
análise é a primeira na ordem dos arquivos de execução, que é a que o juiz da
coleta pontuou, e os vereditos de código de T1, T2 e T3 são computados sobre essa
mesma resposta. Pontuar todas as respostas e tirar a média dos vereditos muda o
veredito de código em \ndupchange\% das células e a média em \ndupmean. No nível do escore,
951 células (11.5\%) têm mais de um escore de juiz, provenientes de reexecuções
em datas diferentes, colapsadas pela \emph{média}: determinística, independente
da ordem de leitura, e usa toda a informação coletada. As quatro políticas
plausíveis (primeira, última, média, aleatória com semente) movem os efeitos em
no máximo 0.06~pp.


\subsection{Acurácia composta por país}

\begin{table}[h]
\centering
\caption{Acurácia composta por país, amostra de 25 países
(tier: GN = Global North, GS = Global South). Ordered by mean.}
\label{tab:s-pais}
\footnotesize
\setlength{\tabcolsep}{4pt}
\begin{tabular}{llr@{\qquad}llr}
\toprule
País & Camada & Média & País & Camada & Média \\
\midrule
Índia & SG & 0.651 & França & NG & 0.534 \\
Austrália & NG & 0.617 & Peru & SG & 0.529 \\
Japão & NG & 0.612 & Colômbia & SG & 0.526 \\
Coreia do Sul & NG & 0.598 & Chile & SG & 0.520 \\
Indonésia & SG & 0.592 & Egito & SG & 0.509 \\
Estados Unidos & NG & 0.587 & Canadá & NG & 0.509 \\
Itália & NG & 0.586 & Quênia & SG & 0.505 \\
África do Sul & SG & 0.575 & México & SG & 0.501 \\
Portugal & NG & 0.562 & Angola & SG & 0.475 \\
Alemanha & NG & 0.557 & Nigéria & SG & 0.461 \\
Reino Unido & NG & 0.554 & Brasil & SG & 0.458 \\
Filipinas & SG & 0.542 & Argentina & SG & 0.433 \\
Bangladesh & SG & 0.542 &  & &  \\
\bottomrule
\end{tabular}
\end{table}

\subsection{Lacuna de camada sob partições alternativas de país}

A divisão Norte/Sul Global segue a classificação da UNCTAD. Uma lacuna visível
apenas sob uma taxonomia seria artefato dela, de modo que a reestimamos sob três
partições baseadas em desenvolvimento dos mesmos 25 países
(Tabela~\ref{tab:s-taxonomia}). A lacuna é significativa sob a UNCTAD e não
significativa sob toda partição traçada pelo IDH, o que é coerente com o gradiente
de IDH nulo relatado no texto principal.

\begin{table}[h]
\centering
\caption{Lacuna de camada sob partições alternativas dos mesmos 25 países.
Teste de permutação com o país como unidade, 10.000 reetiquetagens.}
\label{tab:s-taxonomia}
\footnotesize
\setlength{\tabcolsep}{4pt}
\begin{tabular}{lrrr}
\toprule
Partição & Lacuna (pp) & NG/SG & $p$ \\
\midrule
UNCTAD desenvolvidas/em desenvolvimento (usada no artigo) & $\ntaxunctad$ & 10/15 & $\ntiergapperm$ \\
IDH abaixo da mediana da amostra ($0.781$) & $\ntaxmedtwo$ & 13/12 & $\ntaxmedp$ \\
IDH $<0.800$ (corte ``muito alto'' do PNUD) & $\ntaxvhtwo$ & 12/13 & $\ntaxvhp$ \\
IDH $<0.700$ & $\ntaxlowtwo$ & 20/5 & $\ntaxlowp$ \\
\bottomrule
\end{tabular}
\end{table}

\subsection{Modos representativos de falha, literais}

\begin{table}[h]
\centering
\caption{Quadro 1 --- Modos representativos de falha (literais, levemente truncados).}
\label{tab:s-quadro1}
\footnotesize
\setlength{\tabcolsep}{5pt}
\renewcommand{\arraystretch}{1.15}
\begin{tabular}{@{}>{\raggedright\arraybackslash}p{0.20\linewidth}>{\raggedright\arraybackslash}p{0.74\linewidth}@{}}
\toprule
Modo de falha & Exemplo \\
\midrule
Fabricar um padrão inexistente
& Angola has no national norm. Gemini~2.5~Flash asserted an annual standard of
``25~$\mu$g/m\textsuperscript{3}'' em uma replicação e
``15~$\mu$g/m\textsuperscript{3}'' em outra. O GPT-5 afirmou
``15~$\mu$g/m\textsuperscript{3}'' para a Nigéria, que regula apenas PM\textsubscript{10} (NESREA). Uma resposta fiel se recusaria a declarar um valor que não existe. \\
\addlinespace
Errar o valor vinculante (piso de T1)
& O padrão anual oficial do Brasil é 17~$\mu$g/m\textsuperscript{3} (CONAMA
506/2024). A resposta modal foi 15~$\mu$g/m\textsuperscript{3} (42 respostas), com 10, 20 e 25 também comuns; o valor oficial quase nunca foi devolvido. \\
\addlinespace
Degradação em idioma nativo (H2)
& O padrão anual da Índia é 40~$\mu$g/m\textsuperscript{3}. GPT-OSS~120B
devolveu o valor correto em inglês e resposta \emph{vazia} à mesma pergunta em hindi. \\
\bottomrule
\end{tabular}
\end{table}

% ============================================================
\emph{Modelo misto e reestimação bayesiana.} Um modelo linear misto gaussiano
(MixedLM do \texttt{statsmodels}, REML) com interceptos aleatórios cruzados de país
e modelo dá um efeito do Sul Global de $\beta=\nmixedbeta$ ($\mathrm{EP}=\nmixedse$,
$p=\nmixedp$), marginal em vez de convencionalmente significativo porque cobra do
termo de erro a grande heterogeneidade entre países. Uma reestimação hierárquica
bayesiana com prioris fracamente informativas (\texttt{pymc}; interceptos cruzados
de país e modelo, \texttt{random\_seed=42}) concorda em magnitude e é decisiva
quanto ao sinal: média posterior $\nbayes$, HDI de 94\% $[\nbayeslo,\nbayeshi]$,
$P(\beta<0)=\nbayesp$ ($\hat{R}=\nrhat$).

\emph{E-values.} A lacuna de camada carrega um E-value de $\nevalue$ na estimativa
pontual e $\nevaluelim$ no limite do intervalo mais próximo do nulo (aproximação de
VanderWeele \& Ding 2017 a partir do efeito padronizado); o segundo é o número
operativo, e é um limiar baixo. Nenhum E-value é relatado para o gradiente de
desenvolvimento: seu intervalo já inclui o nulo, de modo que o E-value nesse limite
é $1.00$ e nenhum confundidor é necessário.

\emph{Mediação exploratória (H4).} Um modelo de caminhos padronizado no nível do
país (\texttt{semopy}, $n=25$): $\mathrm{IDH}\rightarrow$ sitelinks $a=+0.67$;
sitelinks $\rightarrow$ acurácia $b=+0.27$; direto $c'=+0.31$; indireto
$ab=+0.18$ com IC bootstrap de 95\% $[-0.27,+0.59]$ (inclui o zero). A cobertura do
país é, portanto, um mediador sugestivo mas não estabelecido, coerente com o
enquadramento do texto principal.

\emph{Checagem da manipulação de persona.} Em 4.349 respostas na condição de
persona, $50.2\%$ reconhecem explicitamente o enquadramento de papel de gestor
público, de modo que o nulo de persona (H6) não é artefato de o enquadramento ser
universalmente ignorado.

% ============================================================
\section{Desvios em relação ao plano de análise pré-especificado}
\label{sec:s-desvios-pt}

O plano de análise exigia que os afastamentos fossem registrados e classificados, e
fixava um teto para eles: o estatuto confirmatório só se mantém se os desvios
afetarem no máximo 15\% das células planejadas. O escopo entregue é de 27.5\% da
contagem planejada de respostas, de modo que o teto é excedido por larga margem, e
o estudo é relatado como exploratório do começo ao fim.

Os desvios principais, com o efeito de cada um sobre a interpretação:

\begin{itemize}\setlength{\itemsep}{3pt}
\item Contagem de respostas. Planejadas 33.600; entregues 8.300 (27.5\%), das
  quais \ncellspt{} entram na análise após remover as células egípcias de T1 sem chave.
  Aciona o próprio teto de 15\% do plano.
\item Registro público. O plano deveria ser depositado em registro público
  antes da coleta confirmatória; foi fixado em documento sob controle de versão e
  nunca depositado. A datação não pode ser verificada de forma independente do
  histórico do repositório, e não reivindicamos pré-especificação pública.
\item Piloto de rubrica. Previsto piloto de 50 prompts com dois
  avaliadores humanos e $\alpha \geq 0.70$ antes de escalar. Não foi executado. A
  confiabilidade é evidenciada a posteriori, em $\alpha = 0.527$, abaixo do limiar
  que o plano fixara.
\item Instrumento de pontuação. O plano previa juiz único. O juiz da coleta não
  foi único: GPT-5-mini pontuou os quinze países pré-especificados e Claude
  Haiku~4.5 os dez da extensão, que contêm sete dos dez países do Norte Global,
  de modo que juiz e tier estavam confundidos nas pontuações julgadas. T1, T2 e
  T3 passam a ser adjudicadas por código contra registros oficiais onde o valor
  devolvido é extraível (99.8\%, 56.3\% e 71.8\% das células); o resíduo e toda
  a T4 são repontuados pela média do painel de três fornecedores; T5 mantém as
  pontuações dos juízes da coleta. Isso muda as conclusões substantivas sobre as
  mesmas respostas: o gradiente cai de $\rho=0.51$ ($p=0.004$) para $\rho=\nrho$
  ($p=\nrhop$); a penalidade de idioma nativo cresce de $-2.1$~pp para $\nnativasigned$~pp e
  se torna significativa nos três idiomas; a razão de chances de T1 passa de
  0.22 sob os juízes para \nort{} sob o código.
\item Chave de T1. A chave do Reino Unido foi corrigida em 21/09/2026 de 10 (meta
  inglesa para 2040) para 20~$\mu$g/m\textsuperscript{3} (valor-limite em vigor
  desde 2020); o Egito não tem valor recuperável e fica fora de T1; Bangladesh foi
  corrigido do diário oficial de 2005 (15) para as Rules de 2022 (35) após
  revisão externa; as escadas oficiais (Brasil, UE, Reino Unido, EUA, Canadá,
  Austrália, Colômbia, Bangladesh) são armazenadas ao lado do valor em vigor; as
  versões do registro carregam um prefixo SHA-256 (\texttt{t1\_registry} b0e5d9e132aab684,
  \texttt{t2\_registry} 703dcd8122f88448, \texttt{t3\_registry} 902de1bc6af893de, \texttt{t4\_reference\_set}
  77bca4679aedc93a). A comparação entre tiers em T1 é reportada com e sem os três países sem
  padrão (\nort{} e \nortexcl) e sob a chave tolerante à escada (\nortladder).
\item Ordenação dos achados. O plano tratava H1 como hipótese primária. H2
  passa a ser relatada como achado principal e o gradiente de H1 como exploratório.
  Reordenação post hoc, motivada pela pontuação corrigida e não por reexame dos
  dados em busca de narrativa melhor; ambas as hipóteses eram pré-especificadas e
  ambas são relatadas.
\item Tradução para idioma nativo. Previstos tradutores humanos nativos com
  retrotradução independente e adjudicação bilíngue. Usou-se tradução por LLM,
  auditada a posteriori por censo de retrotradução dos 90 prompts nativos. O
  artefato contra o qual o plano se protegia está medido e ausente: 90/90 preservam
  o país e a quantidade solicitada.
\item Cobertura de modelos. A cobertura não é uniforme; dois modelos estão
  ausentes de vários países após limitação de taxa e esgotamento de crédito. As
  médias por país são computadas sobre conjuntos de modelos distintos.
\end{itemize}

Uma propriedade adicional do corpus de respostas, que não é desvio do plano mas
consequência de como a coleta foi repetida, pertence a esta seção porque afeta a
reprodutibilidade. De 9.760 chaves únicas (modelo, prompt, replicação), 2.147
(22.0\%) têm mais de uma resposta armazenada, produzidas por reexecuções em datas
diferentes. A política de desempate adotada é a média dos escores, e a análise de
sensibilidade acima mostra que a escolha entre as quatro políticas plausíveis move
os efeitos em no máximo 0.06~pp.

\section{Construção de prompts, template de persona e exemplos verbatim}
% ============================================================

Os prompts foram autorados em inglês para o domínio de política de poluição do ar
com supervisão de domínio (PPGSAU/UTFPR) e validados contra o registro de gabarito
antes de entrar no conjunto de estímulos confirmatório; um subconjunto-piloto
atingiu um limiar inter-avaliadores de $\alpha$ de Krippendorff $\geq 0,70$ antes da
escala. As versões em idioma nativo (português, espanhol, hindi) foram traduzidas e
retrotraduzidas por falantes nativos; qualquer prompt cuja retrotradução produzisse
um desvio semântico foi reescrito ou excluído. Toda afirmação factual em uma
entrada de gabarito precisava resolver para uma URL de fonte oficial ou um DOI
revisado por pares antes de o prompt correspondente ser admitido (o Protocolo de
Verificação de Citações do projeto aplicado já na geração do gabarito, não apenas na
fase de manuscrito).

\paragraph{Manipulação de persona.} A persona é um fator intra-prompt com dois
níveis que mantêm idêntica a pergunta substantiva. O nível \texttt{neutral}
apresenta um pedido simples de informação; o nível \texttt{public\_manager\_env}
prepende um enquadramento fixo de papel. O contraste é mostrado verbatim para o
Brasil, T1:

\begin{quote}\small
\emph{Neutro.} ``What is the current national annual mean ambient air quality
standard for PM2.5 in Brazil? State the regulation that establishes it and whether
it is an intermediate or final target. Limit to 1--3 sentences.''

\smallskip
\emph{Gestor ambiental público.} ``\emph{I am the municipal environmental
management official of S\~ao Paulo, Brazil, and I need to brief decision-makers.}
What is the current national annual mean ambient air quality standard for PM2.5 in
Brazil? State the regulation that establishes it and whether it is an intermediate
or final target. Limit to 1--3 sentences.''
\end{quote}

\noindent O gabarito correspondente no registro é (em tradução): ``Os padrões
nacionais de PM\textsubscript{2.5} do Brasil são definidos pela Resolução CONAMA
506/2024 (que substituiu os artigos pertinentes da Resolução 491/2018). A etapa em
vigor em 2026 é a Meta Intermediária PI-2; a Meta Final adota a diretriz anual da
OMS 2021 de 5~$\mu$g/m\textsuperscript{3}'' (fonte: Resolução CONAMA 506/2024). Os
prompts em idioma nativo aplicam o mesmo template, traduzido; por exemplo, o prompt
T1 em hindi para a Índia é a fiel renderização em devanágari do item neutro em
inglês, reproduzido verbatim no conjunto de prompts liberado. O conjunto completo de
prompts (370 itens confirmatórios: 25 países $\times$ 5 tarefas $\times$ persona,
mais os pares de idioma nativo) é liberado como
\texttt{prompts\_confirmatory.jsonl} com o gabarito, a fonte e a rubrica por item.

% ============================================================
\bibliography{references}

\end{document}
